\section{Equilibrium and Main Results}

This section characterizes equilibrium behavior and derives the main implications of the model. We focus on how stablecoins affect foreign-currency demand, coordination, crisis outcomes, and welfare.

\subsection{Equilibrium Structure}

Given a realization of the public signal $s$, households follow a cutoff strategy. There exists a threshold $x^*(s)$ such that
\begin{equation}
a_i = \begin{cases}
1 & \text{if } x_i \ge x^*(s), \\
0 & \text{otherwise}.
\end{cases}
\end{equation}

Aggregate run demand is therefore
\begin{equation}
Q^{run}(\theta, s) = \bar r \left[1 - \Phi\left(\frac{x^*(s) - \theta}{\sigma_x}\right)\right].
\end{equation}

Total foreign-currency demand is
\begin{equation}
Q^{FX}(\theta, s) = Q^{hedge}(\theta, s) + Q^{run}(\theta, s).
\end{equation}

A crisis occurs if
\begin{equation}
Q^{FX}(\theta, s) \ge R(\theta).
\end{equation}

The cutoff $x^*(s)$ satisfies the indifference condition
\begin{equation}
B - \kappa_r + \int \mathbf{1}\{Q^{FX}(\theta,s) \ge R(\theta)\} L(\theta, \omega_s) f(\theta \mid x^*, s) d\theta = 0.
\end{equation}

\subsection{Information Amplification}

Stablecoins affect equilibrium through the precision of the public signal
\begin{equation}
\sigma_s^2 = \frac{\bar\sigma_s^2}{1 + \rho Q_S}.
\end{equation}

\textbf{Proposition 1 (Information Amplification).} An increase in stablecoin usage increases the weight placed on public information, strengthening coordination across agents.

\textit{Intuition.} Stablecoins generate transparent price signals, improving public information and aligning beliefs.

\subsection{Foreign-Currency Demand}

\textbf{Proposition 2 (Monotonicity).} Total foreign-currency demand and stablecoin usage are increasing in overvaluation $\theta$.

\textit{Intuition.} Higher overvaluation raises the incentive to hold foreign currency, and stablecoins capture an increasing share due to lower costs.

\subsection{Coordination and Crises}

\textbf{Proposition 3 (Coordination Amplification).} Higher stablecoin usage increases the sensitivity of aggregate demand to public signals.

\textbf{Proposition 4 (Crisis Amplification).} Stablecoins increase the probability of crisis by lowering costs and strengthening coordination.

\subsection{Feedback Mechanism}

\textbf{Proposition 5 (Information Feedback).} Stablecoin usage and signal precision are jointly determined and mutually reinforcing.

\subsection{Welfare}

\textbf{Proposition 6 (Non-monotonic Welfare).} Stablecoins increase welfare at low overvaluation and reduce welfare at high overvaluation.

\textit{Intuition.} Efficiency gains dominate at low $\theta$, while coordination-driven crises dominate at high $\theta$.
